\documentclass[10pt,a4paper]{amsart}
\usepackage[margin=19mm]{geometry}
\usepackage{amsmath,amssymb,mathtools}
\usepackage[scr=rsfs,cal=euler]{mathalfa}
\usepackage{xfrac}
\usepackage[unicode,hidelinks,bookmarks=false]{hyperref}
\newcommand{\ie}{i.e.\ }
\newcommand{\cf}{cf.\ }
\newcommand{\resp}{respectively\ }
\newcommand{\Subsection}{\subsection}
\providecommand{\nobreakdash}{\nobreak\hbox{-}}
\setlength{\emergencystretch}{2em}
\multlinegap0pt
\usepackage{amssymb}
\usepackage{mathtools}
\usepackage{xcolor}
\newtheorem{theorem}{Theorem}
\DeclareMathOperator{\B}{\textbf{B}}
\title{Perfectoid Spaces in Multivariate $p$-adic Hodge Theory}
\author{Aprameyo Pal, Rohit Pokhrel}
\date{Source: arXiv:2508.18198v2 (CC BY 4.0)}
\begin{document}
\maketitle

\noindent\emph{Source and licence.} Perfectoid Spaces in Multivariate $p$-adic Hodge Theory. Version arXiv:2508.18198v2 (CC BY 4.0), distributed by arXiv under the Creative Commons Attribution 4.0 International licence, http://creativecommons.org/licenses/by/4.0/, as recorded in the arXiv licence metadata for this exact version and shown on the abstract page https://arxiv.org/abs/2508.18198v2. Full licence notice: \texttt{licenses/arxiv-2508.18198-CC-BY-4.0.txt}.\par\smallskip

\noindent\textit{Editorial scope.} Counted unit: the article's theorem Theorem (source0.tex lines 868--870). The article's complete original proof of it is delivered with it (source0.tex lines 872--918, one \texttt{proof} environment of 47 lines). No mathematics is written, repaired or re-derived: the statement and the proof are the article's own lines, in the article's own notation, and no retained line is substituted or re-flowed.  No further source-local context is needed: every cross-reference of every retained line resolves inside the two delivered spans.  Nothing was omitted from any retained span, so the package carries no omission record.  No retained line contains a citation, so the wrapper carries no bibliography and no bibliography entry.  No retained line contains a figure environment, an image-inclusion command or drawing code, so no image is delivered and the package declares no asset dependency.  Editorial formatting only, in addition: an AMS \texttt{amsart} preamble that restates the article's own declarations and macros used by the retained lines, namely the environments \texttt{theorem} on separate counters, together with the standard packages those lines need.  The retained environments print in this document as Theorem 1 in order of appearance, whereas the article prints the same items with the numbering of its own full document.  The complete native source of the article as distributed in the arXiv e-print for this exact version is bundled unchanged as \texttt{sources/183839/source0.tex}.
\begin{theorem}{\emph{(Open Mapping Theorem)}}\label{Open Mapping Theorem}
	Let $T:E\to F$ be a surjective morphism of Banach spaces. Then $T$ is an open mapping, i.e., if $U$ is open in $E$, then $T(U)\subset F$ is also open.
\end{theorem}

\begin{proof}
	Let $\varpi$ be a pseudo-uniformizer of $K_\Delta$ satisfying $|\varpi|_s=|\varpi|_t$ for all $s,t\in\mathfrak{S}$.  Set $\delta=|\varpi|$ and note that for every $n\in\mathbb{Z}$,
	\begin{align*}
		\B_E(0,\delta^n)&:=\{x\in E\mid \|x-0\|<\delta^n\}\\
		                        &=\{x\in E\mid |\varpi|^{-n}\|x\|<1\}\\
		                        &=\{x\in E\mid \|\varpi^{-n}x\|<1\}\\
		                        &=\varpi^n\B_E(0,1).
	\end{align*}
	Also write $\B_E:=\B_E(0,1)$ which is a subgroup satisfying $\bigcup\limits_{n=1}^\infty\B_E\left(0,\frac{1}{\delta^n}\right)=\bigcup\limits_{n=1}^\infty\varpi^{-n}\B_E=E$. Therefore, we have
	\begin{align*}
		F=\bigcup\limits_{n=1}^\infty \overline{T(\varpi^{-n}\B_E)}=\bigcup\limits_{n=1}^\infty \varpi^{-n}\overline{T(\B_E)}.
	\end{align*}
	By Baire's Category Theorem, there exists $k\in\mathbb{N}^*$ such that $\varpi^{-k}\overline{T(\B_E)}$ has an interior point and hence $\overline{T(\B_E)}$ has an interior point ($x\mapsto \varpi^{-k}x$ is a homeomorphism of $F$). Thus there exists $g\in\B_E$ such that $Tg\in\left(\overline{T(\B_E)}\right)^\circ$, interior of $\overline{T(\B_E)}$. As $\B_E$ is closed under addition, we get
	\begin{align*}
		0\in\left(\overline{T(\B_E-g)}\right)^\circ=\left(\overline{T(\B_E)}\right)^\circ.
	\end{align*}
	Now note that $\{\varpi^n\B_F\}_{n\geq 1}$ form a fundamental system of neighborhoods of $0$ in $F$ because $\lim\limits_{n\to\infty}\varpi^n=0$. There exists $m\geq 1$ such that $\varpi^m\B_F\subset \overline{T(\B_E)}$. By the definition of closure, we have
	\begin{align*}
		h\in\varpi^m\B_F\text{ and }\epsilon>0\implies \text{ there exists } f\in\B_E\text{ such that }\|h-Tf\|<\epsilon.
	\end{align*}
	Let $h$ be arbitrary element of $F$, then there exists $n\in\mathbb{Z}$ such that $|\varpi|^n\leq \|h\|<|\varpi|^{n+1}$ if $n<0$ or $|\varpi|^{n+1}\leq\|h\|<|\varpi|^n$ if $n\geq 0$. Consider
	\begin{align*}
		h'=
		\begin{cases}
			\varpi^{m-n}h &\text{ if } n\geq 0\\
			\varpi^{m-n-1}h &\text{ if } n<0
		\end{cases}.
	\end{align*}
	Note that $h'\in\varpi^n\B_F$ and hence applying above to $h'$, we get:
	\begin{equation}\label{eq 1}
		h\in F \text{ and }\epsilon>0\implies \text{ there exists } f\in
		\begin{cases}
			\varpi^{n-m}\B_E &\text{ if } n\geq 0\\
			\varpi^{n+1-m}\B_E &\text{ if }n<0
		\end{cases}
		\text{ such that }\|h-Tf\|<\epsilon.
	\end{equation}
	Suppose $g\in F$ and $\|g\|<1$. Applying equation (\ref{eq 1}) with $h=g$ and $\epsilon=\delta$ we get:
	\begin{align*}
		\text{ there exists } f_1\in \varpi^m\B_E \text{ such that }\|g-Tf_1\|<\delta.
	\end{align*}
	Again applying equation (\ref{eq 1}) with $h=g-Tf_1$ and $\epsilon=\delta^2$, we get
	\begin{align*}
		\text{there exists }f_2\in\varpi^{m+1}\B_E \text{ such that }\|g-Tf_1-Tf_2\|<\delta^2.
	\end{align*}
	Continuing this process, we can construct a sequence $f_1,f_2,\ldots$ in $E$ such that $f_k\in\varpi^{m+k-1}\B_E$ and $\|g-Tf_1-Tf_2-\cdots-Tf_k\|<\delta^n$. Let $f=\sum\limits_{k=1}^\infty f_k\in \varpi^m\B_E$ which converges because $f_k\to 0$. Indeed, $\varpi^{m+k}\B_E\subset\varpi^m\B_E$, which is a subgroup. Since $T$ is continuous, we obtain $g=Tf$. Therefore, $\B_F\subset T\left (\varpi^{m}\B_E\right)$ which proves the result.
\end{proof}

\end{document}
